\documentclass[11pt]{article}
\usepackage[bottom]{footmisc} 
\renewcommand{\thempfootnote}{\arabic{mpfootnote}}

\usepackage[hidelinks]{hyperref}

\usepackage[utf8]{inputenc}

\usepackage[T1]{fontenc}
\usepackage[ttdefault=true]{AnonymousPro}

\usepackage{multicol}
\setlength{\columnsep}{1cm}

\usepackage{abc}
\usepackage{graphicx}
\usepackage{musicography}

\usepackage[margin=0.5in]{geometry} 
\geometry{letterpaper}


\usepackage{fancyhdr} 
\pagestyle{fancy} % options: empty , plain , fancy
\renewcommand{\headrulewidth}{0pt} % customise the layout...
\lhead{}\chead{}\rhead{}
\lfoot{}\cfoot{\thepage}\rfoot{}


\begin{document}

\thispagestyle{empty}

\begin{abc}[name=whitecockade,options={-O= +c -F single.fmt}]
X:1
T:The White Cockade
M:C|
L:1/8
K:G
|: GA | B2 B2 BAGA | d2 B2 B2 g2 | d2 BB cBAG | FGAB ADGA | BABd cBAG |
B2 d2 g3 a | bagf agfe | d2 B2 B2 :|: Bc | d2 B2 g2 B2 | d3 e d2 Bc |
d2 B2 g2 fg | a2 A2 A2 GA | B2 Bd cBAG | B2 d2 g3 a | bagf agfe | d2 B2 B2 :|
\end{abc}

% \vskip 0.5in
\noindent\makebox[\textwidth][c]{%
  \begin{minipage}{7.0in}
    \setlength{\parindent}{1.5em}
    
This tune is believed to date back to at least 1669 in John Playford’s
tune collection, \textbf{Apollo’s Banquet} (London,
1669).\footnote{John Playford (1623 – 1686/7), London, England, best
  known for his \textbf{The English Dancing Master}
  (1651). \textbf{Apollo’s Banquet for Treble-Violin} (1670) is an
  anthology containing country dances, theater tunes, marches,
  hornpipes, broadside ballads and French marches. This tune is
  apparently marked as `Scots' tune'
  (\url{https://tunearch.org/wiki/Annotation:White\_Cockade\_(1)\_(The)}
  ), but I could NOT figure out which tune it was when I searched the
  scan of the 1687 publication
  (\url{https://imslp.org/wiki/Apollo's\_Banquet\_(Playford,\_John})
  ).  If you find it, let me know which one it is.}  It is later found
in the American publication of \textbf{Evening Amusement}
(Philadelphia, 1796) compiled by Benjamin and Joseph
Carr.\footnote{University of Washington has a digital copy of
  this. \href{https://orbiscascade-washington.primo.exlibrisgroup.com/discovery/fulldisplay?docid=alma99162273569901452&context=L&vid=01ALLIANCE_UW:UW&lang=en&adaptor=Local\%20Search\%20Engine&tab=UW_default&query=any\%2Ccontains\%2C62828387&sortby=rank}{\textbf{Evening
    amusement: containing fifty air's, song's, duett's, dances,
    hornpipe's, reel's, marches, minuett's, \&c, \&c, for 1, and 2
    German flutes or violins. : Price 75 cents.}}  1796; Philadelphia :
  Printed \& sold at B Carr's Musical Repositories Philadelphia and New
  York \& J Carr's Baltimore} The white cockade is thought to either
refer to the white ribbon rosette associated with the Jacobite rebels
in 1715 and 1745 in Scotland and Ireland, or the `bouquets of ribbons
worn by the young women of Munster at weddings and other such
occasions' in the 17 th century.\footnote{tunearch,
  \textit{loc. cit.}}

The White Cockade can be played as a polka\footnote{Scott, Jane Wells. \textbf{Tallahassee Irish Tune Book}, 1996. This is the Tallahassee area Irish Session tune book that I
started with in 1999 and still refer to weekly.} or as a Scottish Measure\footnote{tunearch, \textit{loc. cit.},  See \url{https://socalfolkdance.org/dances/S/Scotch\_Measure.pdf} for Scottish Measure dance steps. YouTube has many videos of people dancing the Scottish Measure.} and can be found in the key of G,
D, or C depending on the publication.

Eloise Hubbard Linscott includes The White Cockade dance steps in her collection, \textbf{Folk Songs of Old
New England} (1939), stating that this dance was being performed in New England in the 1930s.\footnote{Linscott, Eloise Hubbard. \textbf{Folk Songs of Old New England}. Dover, 1962, p 120.}

    {\parindent0pt
      \raggedright
      ~\\
Dance Name: The White Cockade \\
Music Used: The White Cockade \\
Contradance Formation \\
Six to Eight Couples \\
\begin{tabbing}
First lady down the center, first gentleman down the outside and back \qquad \= 8 bars \\
First gentleman down the center, first lady down the outside and back \> 8 bars \\
First couple down the center, back and cast off \> 8 bars \\
Right and left \> 8 bars 
\end{tabbing}
    }
    
A similar White Cockade set dance is being performed here by The Saint Andrew Society of Western
Australia: \url{https://youtu.be/zv1wxtbAIS4}

Vivian and Phil’s Version: \url{https://youtu.be/30hyBXExJmM}

  \end{minipage}
  \vspace*{\fill}
}

\end{document}
